\section{Lexicon}\label{sec:lexicon}

\subsection{The seed lexicon}
The Manual fixes 58 lexical items with their derivations; five numerals (six to ten) and five derived words complete a seed lexicon of 69 entries, listed in full with their nine source forms in Appendix~\ref{app:lexicon}. Its semantic profile follows the design: the vocabulary of the sea and its coast (\emph{meri, kala, räim, sild, tursk, lohi, hülje, kalju, saar, rand, vesi, jää, tuul, lääs, suvi, talv, sügis, päive, boot, dzintar}), the verbs of coastal life (\emph{vidte, kohte, lovte, plavte, žite, žeglte, odpočivte, vite}), the adjectives that describe a sea (\emph{duž, mal, zimne, glembok, solne, pospolit, peleek}), and the function words needed to say something with them. Across the seed lexicon the Hanse Rule decided 30 entries, the Domain Rule 19 and the Family Vote 14 (Table~\ref{tab:rule-stats}); 33 of 63 derivations involved a queue tie-break, so the Rota Rule rotated the queue 33 times before the canonical state was reached.

\subsection{Family balance}
A central design claim is equity between families. Table~\ref{tab:balance} shows which family and which language supplied the winning form for each of the 63 seed derivations. No family dominates: Baltic supplies 20 forms, Slavic 18, Finnic 15 and Germanic 10. Germanic, the family with three members and by far the most speakers, wins \emph{least}, which is the one-family-one-vote rule working as intended; Baltic, with the fewest speakers, wins most, largely through the Hanse Rule (Baltic forms are often the ones shared with Slavic or Germanic) and through Latvian's early place in the queue. Within families the fixed sub-ranking is visible: Estonian decides 13 times to Finnish's 1, Polish 15 to Russian's 3, Latvian 14 to Lithuanian's 6, Swedish 8 to Danish's 2 and German's 1.

\begin{table}[t]
\centering\small
\caption{Winning family and deciding language across the 63 seed derivations.}
\label{tab:balance}
\begin{tabular}{@{}lr@{\quad}lr@{}}
\toprule
Family & Wins & Deciding language & Wins \\
\midrule
Baltic & 20 & Polish & 15 \\
Slavic & 18 & Latvian & 14 \\
Finnic & 15 & Estonian & 13 \\
Germanic & 10 & Swedish & 8 \\
 & & Lithuanian & 6 \\
 & & Russian & 3 \\
 & & Danish & 2 \\
 & & Finnish & 1 \\
 & & German & 1 \\
\bottomrule
\end{tabular}
\end{table}

\subsection{Open vocabulary}
The seed lexicon is a beginning, not a boundary. Any concept can be added by supplying its nine first-dictionary translations, a part of speech and a domain; the procedure does the rest, from the queue state the Manual leaves behind. Table~\ref{tab:newwords} shows six concepts derived this way. \emph{Bread} has no shared skeleton and goes to Baltic by domain (\emph{maiz}); \emph{buy} is a Hanse word between Slavic and Germanic (\emph{kup}, from Polish \emph{kupić}, Swedish \emph{köpa}, German \emph{kaufen}); \emph{child} unites Latvian \emph{bērns} with Scandinavian \emph{barn}; \emph{market} is the Hanseatic word par excellence, T-R-K in Finnic (\emph{turg}), Baltic (\emph{tirgus}), Slavic (\emph{targ}) and Germanic (\emph{torg}); \emph{ship} joins Finnic \emph{laiva} with Baltic \emph{laivas}; \emph{hand} joins Baltic \emph{roka} with Slavic \emph{ruka}. These are the two faces of the lexicon: where the shore already shares a word, Baltmeri inherits it; where it does not, the domain's family speaks.

\begin{table}[t]
\centering\small
\caption{Six concepts derived beyond the Manual, from the canonical queue.}
\label{tab:newwords}
\setlength{\tabcolsep}{3pt}\scriptsize
\begin{tabular}{@{}llp{1.5cm}p{1.4cm}p{2.7cm}@{}}
\toprule
Concept & Baltmeri & Rule & Decided by & Shared skeleton \\
\midrule
bread & maiz & Domain (farm) & Latvian & --- \\
buy & kup- & Hanse & Polish (tie) & K-P: pl ru sv da de \\
child & barn & Hanse & Danish (tie) & P-R-N: lv sv da \\
market & torg & Hanse & Russian (tie) & T-R-K: et lv lt pl ru sv \\
ship & laiv & Hanse & Lithuanian (tie) & L-V: fi et lt \\
hand & rok & Hanse & Latvian (tie) & R-K: lv ru \\
\bottomrule
\end{tabular}
\end{table}

\subsection{Derivation order and the coining of words}\label{sec:eval-queue}
Because the queue rotates on every tie, the form of a new word depends on the state of the queue when it is coined, and hence on what was coined before it. The Manual resolves this for its own vocabulary by fixing the order of derivation. For open vocabulary we adopt two conventions. First, the queue state after the whole Manual has been derived---the \emph{canonical queue}, which is
\begin{quote}
Finnish, German, Russian, Polish, Danish, Swedish, Lithuanian, Estonian, Latvian
\end{quote}
---is the starting state for every new derivation; within a single text, new concepts are derived in order of first appearance. Second, a concept's first derivation is recorded and reused forever after (``first derivation wins''), so that a word, once coined, never changes, while two speakers coining different words cannot perturb each other's tie-breaks. The canonical queue is itself a small monument to the Rota Rule: Swedish, first in the initial queue, has been rotated to sixth, and Latvian, which decided fourteen ties, is last.
